\section{Cobertura medios de comunicación}


\subsection{Presentación del caso}

Una compañía quiere lanzar un nuevo producto al mercado, para lo que realizará una campaña publicitaria para dar a conocer el nuevo producto. Tras realizar un estudio de mercado, la compañía ha identificado varios segmentos de mercado (audiencias) interesados en el nuevo producto. 

A su vez existen diferentes medios de comunicación para dirigirnos a los diferentes segmentos de mercado, pero ha comprobado que, de los medios de comunicación existentes, ninguno alcanza a todos los segmentos de mercado, por lo que resulta necesario seleccionar varios medios a la vez para llegar a todas las audiencias. 

El estudio de mercado ha identificado seis tipos de audiencias potenciales diferentes y la posibilidad de seleccionar ocho medios de comunicación, con las incidencias y costes asociados que se muestran en la tabla correspondiente.  

Se trata de decidir qué medios de comunicación deben seleccionarse para alcanzar todos los segmentos de mercado con el menor coste posible.

\subsubsection{Modelo base}

\paragraph{Variables}\mbox{}\\

\begin{tabular}{l c l}
$x_j$ & : & variable binaria que vale 1 si se selecciona el medio $j$, 0 en caso contrario, para $j = 1, \dots, 8$. \\
\end{tabular}
\\
\\

\paragraph{Parámetros}\mbox{}\\

\begin{tabular}{l c l}
$C_j$ & : & coste del medio $j$ \\
$A_{ij}$ & : & parámetro que vale 1 si el medio $j$ alcanza a la audiencia $i$, 0 en caso contrario, para $i = 1, \dots, 6$.\\
\end{tabular}
\\
\\

\paragraph{Restricciones}\mbox{}\\

Cada audiencia debe ser alcanzada por al menos un medio de comunicación seleccionado:
\begin{equation}
\sum_{j=1}^{8} A_{ij}x_j \geq 1 \qquad \forall i = 1,\dots,6
\end{equation}
\\
\\

\paragraph{Función objetivo}\mbox{}\\

\begin{equation}
\mbox{min. } z = \sum_{j=1}^{8} C_j x_j
\end{equation}
\\
\\

\paragraph{Condiciones de no negatividad y naturaleza}\mbox{}\\
\begin{equation}
x_j \in \{0,1\} \qquad \forall j = 1,\dots,8
\end{equation}
\\
\\

\subsubsection{Obligación de elegir al menos un medio de TV}

Si se dispone de dos medios de televisión, TV diurna y TV nocturna, y debemos seleccionar al menos uno de ellos, la restricción correspondiente es:

\begin{equation}
x_{\text{TVd}} + x_{\text{TVn}} \geq 1
\end{equation}

\noindent
De este modo, el modelo garantiza que la campaña incluirá al menos uno de los dos medios televisivos.

\subsubsection{Obligación de elegir al menos un medio de TV si se elige algún tipo de cartel}

Si existen dos medios de carteles (por ejemplo, cartel urbano y cartel carretera), y queremos imponer que si se elige alguno de ellos, debe elegirse al menos una TV, la restricción lógica es:

\begin{align}   
x_{\text{TVd}} + x_{\text{TVn}} \geq x_{\text{CU}} \\
x_{\text{TVd}} + x_{\text{TVn}} \geq x_{\text{CC}}
\end{align}
\noindent
Esto asegura que, si se escoge algún cartel ($x_{\text{CU}}=1$ o $x_{\text{CC}}=1$), el modelo deberá seleccionar al menos una TV.

\subsubsection{Obligación de elegir al menos un tipo de cartel si se selecciona alguna TV}

En este caso la implicación es la inversa: si se escoge al menos una de las dos televisiones, debe seleccionarse al menos uno de los dos tipos de carteles.  
La restricción sería:



\begin{align}   
x_{\text{CC}} + x_{\text{CU}} \geq x_{\text{TVd}} \\
x_{\text{CC}} + x_{\text{CU}} \geq x_{\text{TVn}}
\end{align}
% \begin{equation}
% x_{\text{CU}} + x_{\text{CC}} \geq x_{\text{TVd}} + x_{\text{TVn}} - 1
% \end{equation}

\noindent
Esta ecuación impide que haya selección de TV sin carteles, pero no obliga a más de uno.

\subsubsection{Incompatibilidad entre TV diurna y carteles}

Si la selección de TV diurna impide contratar cualquiera de los dos tipos de carteles, debemos expresar la incompatibilidad mediante las siguientes restricciones:

\begin{align}
x_{\text{TVd}} + x_{\text{CU}} &\leq 1 \\
x_{\text{TVd}} + x_{\text{CC}} &\leq 1
\end{align}

\noindent
Así, si $x_{\text{TVd}}=1$, ambos $x_{\text{CU}}$ y $x_{\text{CC}}$ deben ser 0.

\subsubsection{Obligación de seleccionar el periódico financiero si se eligen  revistas y TV nocturna}

Si $x_{\text{R}}$ representa la selección de revistas y $x_{\text{TVn}}$ la de TV nocturna, y $x_{\text{PF}}$ el periódico financiero, la condición de que si se eligen ambas (revistas y TV nocturna) se deba elegir también el periódico financiero se expresa así:

\begin{equation}
x_{\text{PF}} \geq x_{\text{R}} + x_{\text{TVn}} - 1
\end{equation}

\noindent
Esto garantiza que, si $x_{\text{R}}=1$ y $x_{\text{TVn}}=1$, entonces $x_{\text{PF}}=1$.



%%%%%%%%%%%%%%%%%%%%%%%%%%%%%%%%%%%%%%%%%%%%%%%%%%%%%%%%%%%%%%%%%%%%%%%%%%%%%%%%%%%%%%%%%%%%%%%%%%%%%%%%%
% VARIANTE 6
%%%%%%%%%%%%%%%%%%%%%%%%%%%%%%%%%%%%%%%%%%%%%%%%%%%%%%%%%%%%%%%%%%%%%%%%%%%%%%%%%%%%%%%%%%%%%%%%%%%%%%%%%
\newpage
\subsubsection{Al menos tres audiencias deben ser alcanzadas más de una vez}

Para imponer que, además de cubrir todas las audiencias, al menos tres de ellas sean alcanzadas por más de un medio, se definen variables binarias auxiliares:



\begin{tabular}{l l}
$y_i$ : & variable binaria que vale que valen 1 si la audiencia $i$ es alcanzada más de una vez, \\
& 0 en caso contrario, para $i = 1, \dots, 6$. \\
\end{tabular}


\begin{align}
\sum_{j=1}^{8} A_{ij}x_j - 1 &\leq 6\times y_i \qquad \forall i=1,\dots,6 
\label{eq:cobertura6-1}
\end{align}

La restricción \ref{eq:cobertura6-1} impone que si el número de medios utilizados es 2 o mayor, entonces se cumple que $\sum_{j=1}^{8} A_{ij}x_j - 1 \geq 1$ y la forma de que la restricción se cumpla es que $y_i=1$


\begin{align}
y_i \leq & \sum_{j=1}^{8} A_{ij}x_j - 1 \qquad \forall i=1,\dots,6 
\label{eq:cobertura6-2}
\end{align}

Por su parte, la restricción \ref{eq:cobertura6-2} impone que si solo un medio incide sobre una audiencia, es decir, si $\sum_{j=1}^{8} A_{ij}x_j=1$, la parte de la izquierda será 0 y, por lo tanto, $y_i\leq 0$, es decir, $y_i\leq 0$, con lo que la variable $y_i$ tomará el valor correspondiente cuando la audiencia $i$ no recibe la publicidad por más de un medio, es decir, 0.

En resumen, entre las restricciones \ref{eq:cobertura6-1} y \ref{eq:cobertura6-2} se asegura que si se indice sobre la audiencia $i$ solo una vez $y_i=0$ y si se incide dos o más veces $y_i=1$


\begin{align}
\sum_{i=1}^{6} y_i &\geq 3
\label{eq:cobertura6-3}
\end{align}

\noindent
Por último, la restricción \ref{eq:cobertura6-3} fuerza a que al menos tres audiencias tengan cobertura múltiple.
